\chapter{Basic Notions} 
\label{chap:notions}

To make this work more self-contained, we briefly introduce basic notions and concepts used in the later chapters. 
Details can be found in \update{doplnit odkazy na nejaka skripta, klidne nekolik; jednim z nich by mel byt Hartley, Zisserman: Multiple View Geometry}{\cite{multipleview}}.
% TODO v tehle sekci ale chceme rozebirat jen matematicke a image processing pojmy, ne ty softwarove, coz je otazka, jestli bychom to nemeli nejak zminit... ale mozna ani ne... nevim

\begin{framed} 
% TODO 
\todo{Co budeme potrebovat zavest:} 
\update{nasledujici podkapitolky podle tohoto seznamu}{} 
\begin{itemize} 
\item convolution 
\item image derivatives 
\item gaussian image derivatives 
\item sum of absolute differences 
\item entropy, mutual information 
\item hessian, laplacian 
\item homogeneous coordinates 
\item camera model (mozna do sekce o MVG?) 
\item fundamental matrix (mozna do sekce o MVG?) 
\end{itemize} 
\end{framed} 

\section{Convolution}

\term{Convolution} is an operation on two functions used in signal processing in general.
In image processing \term{convolution} performs applying some filter (kernel) over an image.
The output of such \term{convolution} can be blurred or sharpened image for example.

\begin{definition}
\term{The convolution of functions $f$ and $g$} is an operation defined as:
\begin{align*}
(f\convolution g)(t)&=\inftyint f(\tau)g(t -\tau) d\tau.
\end{align*}
In image processing $f$ would represent the image function and $g$ would be the filter applied on it.
\end{definition} 

\section{Image derivatives}
\label{sec:imder}
\term{Image derivatives} are local derivatives of an image function.
They are proper way how to express changes in an image.
By a change in an image we assume an acute change of a pixel intensity that indicates presence of an edge.
After derivation the change transforms to maxima of the image function.
That is one of the methods how edges can be detected.
In discrete space we can approximate derivatives by local differences as:
\begin{align*}
\imd{ \partial x}(x, y)& = I(x + 1, y) - I(x-1, y)\\
\imd{ \partial y}(x, y)& = I(x, y + 1) - I(x, y - 1),
\end{align*}
where $I$ is the input image and \vect{(x, y)} are coordinates of the pixel.

\section{Gaussian image derivatives}

\update{}{k teto podkapitolce jsem cerpala z nejake prezentace, co jsem nasla na internetu. nejsem si jista, jestli jsem to spravne pochopila.}

In image processing and \cv\ scale-space is a set of down-sampled and smoothed images.
Formally, it is a group of derived signals $\L$ defined as \term{convoltion} of image function $\func$ and Gaussian kernel $\gauss$
%\begin{align*}
%\gauss = \frac{1}{2\pi t} e^{-(x^2+y^2)/2t}
%\end{align*}
so that
\begin{align*}
\L = \gauss \convolution \func,
\end{align*}
where $t$ is the particular scale.

To this scale-space representation we can apply local derivatives at any scale.
Equivalently, scale-space derivatives can be computed by convolving the original image $f$ with Gaussian derivative operators which are derivatives of the Gaussian function.
For this reason they are often also referred to as \term{Gaussian derivatives}.

\section{Sum of absolute differences}

\term{Sum of absolute differences} describes the difference of intensities of particular areas of two images.
In the context of image processing it is used to measure disparity of the areas and by minimizing the difference to find the best matching points.
By computation of \term{sum of absolute differences} we can estimate the disparity map.

\section{Entropy}

In computer science, \term{entropy} (sometimes also called \term{Shannon entropy}) is a measurement of the uncertainty in the image values.
\term{Entropy} can be also perceived as a measure of histogram dispersion.


\begin{definition}
The entropy $H$ of a random variable $X$ and a probability $P(X)$ is defined as
\begin{align*}
  H(X) = E[-\ln(P(X))] = - \sum_k{P(k)\log P(k)},
\end{align*}
where $E$ is expected value operator and $k$ is a possible value of the $X$ variable.
\end{definition}

\section{Mutual information}

\term{Mutual information} measures the mutual dependence of two variables and says how much one random variable tells us about another.
\begin{definition}
\term{The mutual Information} of two variables $I_1$ and $I_2$ is a measure of uncertainty (entropy $H$) in a probability density function:
\begin{align*}
\MI{I_1, I_2} = \H{I_1} + \H{I_2} - \H{I_1, I_2}.
\end{align*}  
\end{definition}
We can also perceive \term{mutual information} as a reduction in uncertainty about one random variable given knowledge of another.
High mutual information indicates a large reduction in uncertainty; vice versa low mutual information indicates a small reduction.
In the context of image processing high mutual information demonstrate high entropy and a white noise in an image.


\section{Laplacian}

In section \ref{sec:imder} we mentioned \term{image derivatives} that are useful for detection of high variation of intensity.
After derivation they occurred as maxima of the image function.
If we go even further and take the second derivative we realise that the intensity changes transform to zero.
Considering n-dimensional space of the function, we get $n$ second partial derivatives.
If we sum them we get the \term{Laplacian}.

\begin{definition}
The Laplacian in n-dimensional space is the divergence of a gradient of function $f$:
\begin{align*}
\nabla ^{2}f&= \sum_{i=0}^{n} \fpp{ \partial x_i^2 }.
\end{align*}
In image processing we usually consider 2-dimensional space, thus the \term{Laplacian} is
\begin{align*}
\nabla ^{2}f(x, y)&=\fpp{\partial x^2} + \fpp{\partial y^2}.
\end{align*}
\end{definition} 

\section{Homogeneous coordinates}
In Euclidean coordinate system infinity does not exist.
However, many geometric concepts are simplified when infinity is included.
An example of such a form of geometry that includes infinity is \term{projective geometry}.

Suppose we have a point $(x, y)$ in Euclidean plane.
This point in projective geometry is expressed by extending it to the triplet value of $(x, y, 1)$.
Generally, a n-dimensional point in Euclidean space is represented as a point in a (n+1)-dimensional projective space.
The point $(kx, ky, k)$ is certainly the same point as $(x, y, 1)$ for any $k \neq 0$.
These are called \term{homogenous coordinates} of the point.
To get back the Euclidean coordinates we simple divide $x$ and $y$ by $k$.
We can notice that any of $(x, y)$ from Euclidean space corresponds to the point $(x, y, 0)$, because the operations $x/0$ and $y/0$ are undefined.
This is how we get the point in infinity; points in infinity are those points which have zero as the last coordinate.


\section{Projection camera model}

\term{Projection camera model} represents a model of central perspective projection.
It maps points from 3D (world) to 2D (image).

In Figure \ref{fig:camera} we can see the camera geometry.
The \term{camera centre} is placed at the coordinate origin.
We should notice that the image plane is in front of camera.
The line perpendicular to the \term{image plane} that origins in the centre is called \term{principal axis} of the camera and its intersection with the image plane is \term{principal point}.

%\begin{align*}
%\begin{bmatrix}
 % x_1\\
  %x_2\\
  %x_3
%\end{bmatrix}&= 
%\T
%\begin{bmatrix}
 % X_1\\
  %X_2\\
  %X_3\\
  %X_4
%\end{bmatrix},
%\end{align*}
%where $\T$ is the transformation from homogenous coordinates to Euclidean space.
\begin{figure}[h]
  \label{fig:camera}
  \centering{\includegraphics[width=96mm]{img/perspectivecamera.png}}
  \caption{\todo{tento obrazek jsem stahla z googlu, mel by se nahradit}}
\end{figure}

\section{Fundamental matrix}

In \cv\ \term{fundamental matrix} represents the relation of points between two stereo images.
Suppose we have two stereo images and we choose a point $\x$ in the first one.
The knowledge of its position constraints the position of the corresponding one in the second image.
The corresponding point denoted by $\xd$ lies on the line defined by $\F\x$ where $\F$ is the fundamental matrix.
This line is called \term{epipolar line}.

\begin{definition}
The \term{fundametal matrix $\F$} for a pair of stereo images is a $3 \times 3$ matrix which satisfies 
\begin{align*}
\xd^T \F \x = 0
\end{align*}
for corresponding points $\x$ and $\xd$.
\end{definition}

\section{Hessian matrix}

% feature detectors and interest points descriptors employ the Hessian matrix of the image.
The Hessian matrix describes a second-order behaviour of a function around a particular point. 

\begin{definition} 
\term{The Hessian matrix of a function $f: \Rn \to \R$ at $\x \in \Rn$} is a matrix $H_f(\x)$ of the second order partial derivatives of $f$ evaluated at $\x$: %~=~(x_1, \ldots, x_n)$: 
\begin{align*} 
H_f(\x) := 
\begin{pmatrix} 
\fpp{ \partial x_1^2 }              &   \fpp{ \partial x_1 \partial x_2 }   &   \ldots   &   \fpp{ \partial x_1 \partial x_n }   \\ 
\fpp{ \partial x_2 \partial x_1 }   &   \fpp{ \partial x_2^2 }              &   \ldots   &   \fpp{ \partial x_2 \partial x_n }   \\ 
\hdotsfor[2]{4} \\ 
\fpp{ \partial x_n \partial x_1 }   &   \fpp{ \partial x_n \partial x_2 }   &   \ldots   &   \fpp{ \partial x_n^2 }  
\end{pmatrix}. 
\end{align*} 
If any of the partial derivatives on the right-hand side are undefined, we say that the Hessian matrix is also undefined.
% \[ H(f, \x) = \left(\begin{array}{c}  
% \frac{\partial^2 f}{\partial x_1^2} \frac{\partial^2 f}{\partial x_1 \partial x_2} 
% \cdots \frac{\partial^2 f}{\partial x_1 \partial x_n}\\
% \frac{\partial^2 f}{\partial x_2 \partial x_1} \frac{\partial^2 f}{\partial x_2^2} 
% \cdots \frac{\partial^2 f}{\partial x_2\partial x_n}\\
% \cdots \cdots \cdots \\
% \frac{\partial^2 f}{\partial x_n \partial x_1} \frac{\partial^2 f}{\partial x_n \partial x_2} 
% \cdots \frac{\partial^2 f}{\partial x_n^2}
% 
% \end{array}\right)  \]
\end{definition} 
Within the context of image processing, the function $f$ typically corresponds to the input image and derivatives are replaced either by differences between the intensity levels of neighbouring pixels or Gaussian derivatives. 

The Hessian matrix forms a basis for a basic feature detector (called Hessian detector), which selects the image positions $\x$ locally maximizing $\det(H(I, \x)),$ where $I$ is the input image.
This leads to a detection of corners and features forming highly textured areas of size roughly comparable to the variance of the Gaussian kernel used to compute the Gaussian derivatives. % TODO zkontrolovat, jestli to neni nesmysl

\section{Integral images}

% Later in this work we work with the term of integral images. 
\term{An integral image}, also known as \term{summed area table}, allows fast and efficient computation of a sum of image intesity values inside an arbitrary rectangular area.
A pixel of an integral image represents the sum of all of the original image's pixels that lie to the left and above the considered position: 
\begin{equation*}
\K_I(\x) := \sum_{i \le x} \sum_{j \le y} I(i,j),
\end{equation*}
where $\K$ is the resulting integral image, \emph{I} is the input image image, and $\x = \vect{(x, y)}^{T}$ is a location of a pixel.

An advantage of the integral image is that we are able to compute it using only one pass through the original image. 
Moreover, once we have calculated the integral image, only three integer operations and four memory accesses are required to calculate the sum 
of the original intensities inside any rectangular region (see Figure \ref{fig:integral}).

\begin{figure}[h]
  \label{fig:integral}
  \centering{\includegraphics[width=96mm]{img/integral_image.png}}
  \caption{The sum of any rectangular region can be calculated by only three additions: \update{napsat konkretni vzorec}{$\sum = C - D - B + A$}.}
 \end{figure}
